\clearpage
\appendix

\section{Toàn bộ các đợt vượt Q95}
\label{app:events}

Bảng~\ref{tab:all-events} liệt kê đầy đủ 15 đợt, bao gồm ngày bắt đầu, ngày kết
thúc, số ngày, ngày đỉnh và lưu lượng đỉnh. Các đợt được sắp xếp theo thời gian.

\begin{table}[H]
\centering
\footnotesize
\setlength{\tabcolsep}{3.2pt}
\begin{tabular}{llrrrrr}
\toprule
Mã & Bắt đầu & Kết thúc & Số ngày & Ngày đỉnh & Đỉnh (m³/s) & Ngày vượt \\
\midrule
E01 & 02/06/2017 & 02/06/2017 & 1 & 02/06/2017 & 419{,}44 & 1 \\
E02 & 04/06/2017 & 07/06/2017 & 4 & 05/06/2017 & 422{,}87 & 4 \\
E03 & 10/06/2017 & 10/06/2017 & 1 & 10/06/2017 & 418{,}28 & 1 \\
E04 & 13/06/2017 & 28/08/2017 & 77 & 06/07/2017 & 453{,}52 & 77 \\
E05 & 02/09/2017 & 05/09/2017 & 4 & 03/09/2017 & 423{,}71 & 4 \\
E06 & 15/09/2017 & 14/11/2017 & 61 & 11/10/2017 & 439{,}97 & 61 \\
E07 & 09/01/2018 & 09/01/2018 & 1 & 09/01/2018 & 417{,}66 & 1 \\
E08 & 13/01/2018 & 16/01/2018 & 4 & 14/01/2018 & 423{,}21 & 4 \\
E09 & 19/01/2018 & 19/01/2018 & 1 & 19/01/2018 & 417{,}91 & 1 \\
E10 & 21/01/2018 & 23/01/2018 & 3 & 23/01/2018 & 418{,}78 & 3 \\
E11 & 04/02/2018 & 07/02/2018 & 4 & 06/02/2018 & 420{,}85 & 4 \\
E12 & 11/03/2018 & 13/03/2018 & 3 & 12/03/2018 & 418{,}79 & 3 \\
E13 & 29/11/2019 & 29/11/2019 & 1 & 29/11/2019 & 417{,}40 & 1 \\
E14 & 02/02/2020 & 02/02/2020 & 1 & 02/02/2020 & 417{,}46 & 1 \\
E15 & 04/02/2020 & 03/05/2020 & 90 & 10/03/2020 & 457{,}49 & 90 \\
\bottomrule
\end{tabular}

\caption{Toàn bộ các đợt lưu lượng vượt Q95 trong giai đoạn 2010--2023.}
\label{tab:all-events}
\end{table}

\section{Thống kê lưu lượng theo năm}
\label{app:yearly}

Bảng~\ref{tab:yearly-stats} bổ sung thống kê theo năm cho phân tích ở câu hỏi 1.
Các giá trị được tính trên toàn bộ ngày trong năm, không phải trung bình của các
tháng.

\begin{table}[H]
\centering
\footnotesize
\setlength{\tabcolsep}{3.2pt}
\begin{tabular}{rrrrrrrr}
\toprule
Năm & n & Mean & Median & SD & Min & P95 & Max \\
\midrule
2010 & 365 & 129{,}01 & 128{,}44 & 15{,}22 & 103{,}28 & 154{,}22 & 160{,}82 \\
2011 & 365 & 139{,}65 & 135{,}27 & 19{,}41 & 111{,}01 & 172{,}77 & 185{,}86 \\
2012 & 366 & 176{,}56 & 193{,}78 & 45{,}89 & 94{,}39 & 240{,}40 & 246{,}60 \\
2013 & 365 & 277{,}89 & 278{,}14 & 15{,}51 & 236{,}31 & 305{,}51 & 328{,}24 \\
2014 & 365 & 358{,}01 & 356{,}51 & 25{,}17 & 316{,}27 & 400{,}14 & 412{,}01 \\
2015 & 365 & 371{,}48 & 370{,}66 & 9{,}29 & 350{,}52 & 388{,}11 & 402{,}06 \\
2016 & 366 & 373{,}73 & 370{,}38 & 13{,}17 & 346{,}57 & 394{,}87 & 398{,}21 \\
2017 & 365 & 411{,}57 & 408{,}96 & 20{,}68 & 373{,}61 & 448{,}42 & 453{,}52 \\
2018 & 365 & 377{,}56 & 377{,}62 & 25{,}01 & 328{,}91 & 416{,}53 & 423{,}21 \\
2019 & 365 & 391{,}50 & 390{,}80 & 12{,}99 & 362{,}92 & 412{,}08 & 417{,}40 \\
2020 & 366 & 382{,}61 & 366{,}30 & 36{,}49 & 330{,}80 & 443{,}42 & 457{,}49 \\
2021 & 365 & 305{,}25 & 313{,}55 & 32{,}81 & 244{,}42 & 342{,}80 & 351{,}85 \\
2022 & 365 & 231{,}00 & 230{,}35 & 9{,}59 & 209{,}54 & 246{,}67 & 251{,}39 \\
2023 & 365 & 213{,}61 & 212{,}16 & 16{,}43 & 180{,}26 & 239{,}48 & 243{,}57 \\
\bottomrule
\end{tabular}

\caption{Thống kê lưu lượng theo năm, 2010--2023.}
\label{tab:yearly-stats}
\end{table}

\section{Manifest và kiểm tra tính nhất quán}
\label{app:manifest}

Bảng~\ref{tab:manifest} tóm tắt các thông tin nhận dạng của lần chạy. Các thông tin
này được lưu trong \texttt{run\_manifest.json}, \texttt{data\_quality.json} và
\texttt{model\_selection.json}. Nếu một trong các giá trị này thay đổi, kết quả
trong báo cáo không còn tương ứng với lần chạy mới.

\begin{table}[H]
\centering
\small
\begin{tabularx}{\linewidth}{@{}lX@{}}
\toprule
\textbf{Trường} & \textbf{Giá trị} \\
\midrule
Run ID & \texttt{\RunId} \\
Dataset hash & \texttt{f83d984747fe\ldots} (đầy đủ trong manifest) \\
Config hash & \texttt{b45e27e26f42\ldots} (đầy đủ trong manifest) \\
Seed & 42 \\
Số ngày dữ liệu & \DatasetN \\
Số gốc Test & 724 \\
Số cặp Test mỗi mô hình & 5.068 \\
Ngưỡng Q95 & \QHigh{} m\textsuperscript{3}/s \\
Số đợt vượt Q95 & \NEvents \\
Số bảng / hình full & 14 / 14 \\
Python & 3.12.3 \\
pandas & 3.1.0 \\
NumPy & 2.5.4 \\
statsmodels & 0.15.0 \\
LightGBM & 4.7.0 \\
Matplotlib & 3.11.2 \\
\bottomrule
\end{tabularx}
\caption{Thông tin nhận dạng và phiên bản phần mềm của lần chạy.}
\label{tab:manifest}
\end{table}

Các kiểm tra tự động đã chạy gồm: audit dữ liệu, kiểm tra split không xáo trộn,
kiểm tra biên ngày mục tiêu, 15 bài test Python và build website. Kết quả test
không thay thế kiểm tra khoa học, nhưng bảo đảm pipeline không bị lỗi cú pháp hoặc
lỗi dữ liệu cơ bản.

\section{Hướng dẫn tái lập}
\label{app:reproduce}

Các lệnh dưới đây chạy từ thư mục gốc của dự án. Môi trường được tạo bằng
\texttt{uv}; có thể dùng \texttt{pip} nếu đã có Python 3.12.

\begin{verbatim}
cd code
uv venv --system-site-packages .venv
uv pip install --python .venv/bin/python -e .
.venv/bin/python -m colorado_flow.pipeline --stage all
\end{verbatim}

Lệnh trên chạy lần lượt audit, EDA, features, train, freeze, evaluate và export.
Kết quả nằm trong \texttt{code/results}; bảng và hình cho báo cáo nằm trong
\texttt{latex/} và \texttt{latex\_full/}.

Để build website:

\begin{verbatim}
cd code/web
npm ci
npm run build
\end{verbatim}

Để build báo cáo đầy đủ bằng Tectonic:

\begin{verbatim}
cd latex_full
tectonic -X compile main.tex
\end{verbatim}

Nếu dùng Overleaf, upload thư mục \texttt{latex\_full} và đặt \texttt{main.tex}
làm tệp chính. Các tệp bảng, hình và số liệu đã được sinh sẵn; không cần chạy lại
Python trên Overleaf.

\section{Ghi chú về hai phiên bản báo cáo}
\label{app:versions}

Báo cáo này là bản đầy đủ, được mở rộng từ bản nộp 5 trang theo yêu cầu của người
học. Bản 5 trang được giữ trong \texttt{latex/} và
\texttt{deliverables/} (tệp PDF cùng mã sinh viên); bản đầy đủ được build
từ \texttt{latex\_full/main.tex}. Hai phiên bản dùng chung dữ liệu, cùng run ID và
cùng bộ kết quả, nên không có mâu thuẫn số liệu giữa bản ngắn và bản dài.

Việc giữ cả hai phiên bản giúp người đọc chọn mức độ chi tiết phù hợp: bản ngắn để
nộp đúng giới hạn trang, bản dài để trình bày đầy đủ phương pháp, chẩn đoán và
thảo luận. Khi có thay đổi dữ liệu hoặc mô hình, cả hai cần được build lại từ cùng
một lần chạy pipeline.
